\chapter{\textbf{Literature review and Project Contribution}}

Cable-suspended cooperative payload transportation using UAVs has
garnered considerable attention in recent years. Researchers have
explored a wide range of techniques, from classical methods such as
Optimal Control [6] and closed-loop approaches with carrier trajectories 
calculated online using a Nonlinear Programming (NLP) solver, to 
learning-based methods such as Imitation Learning and Reinforcement Learning 
[7] [8] [9], with each offering its own advantages and disadvantages. 

\section{Multi-rotor cooperative transportation}

In [6], the authors propose a centralized kinodynamic
motion planner that formulates a finite-horizon Optimal Control Problem
(OCP) and solves it online in a receding-horizon fashion. The approach
aims to enable high-speed, high-acceleration manipulation of a
cable-suspended payload while accounting for the strong dynamic coupling
between the UAVs, cables, and payload. While it achieves strong results,
it relies on a centralized kinodynamic planner, which may pose challenges
for fully onboard deployment of the complete control system. 

In [7], a
centralized kinodynamic planner is distilled using DAgger-based imitation
learning into a homogeneous, decentralized UAV policy\footnote{\textbf{Distillation} is a form of imitation learning that involves transferring knowledge from a teacher policy, typically obtained through planning or expert demonstration, to a smaller or more deployable student policy by training the latter to reproduce the teacher's actions. \textbf{DAgger (Dataset Aggregation)} is an imitation-learning algorithm that iteratively collects expert actions on states visited by the learned policy and aggregates these samples into the training dataset, thereby reducing the distribution-shift problem encountered when the learner visits states not represented in the original demonstrations.}. This allows the
approach to achieve agile full-pose cooperative manipulation using only
local observations and without communication. While the policy was
decentralized, it was not deployed onboard during the actual testing due
to its high computational cost. 

In [8], a decentralized cooperative
control policy was learnt directly through multi-agent reinforcement
learning (MARL), trading explicit trajectory-planning capability for much
cheaper computational cost and greater robustness to heterogeneous agents
and agent failure. Additionally, the policy was ony tasked to learn high level
commands/action such as reference acceleration. The action is then handed to 
a low level controller that executes the actual tracking. In this case, 
since the policy had dramatically lower computational cost, actual 
onboard execution was possible.

In [9], instead of relying on a low level controller, an end-to-end 
decentralized cooperative control policy was learnt. The policy takes as
input the sensor measurements and outputs the PWM directly to drones' motors.
This eliminates the need for a low level controller and simplifies the 
deployment. However, while the policy did not rely on explicit communication
between the drones, it relied on inter-drone measurements, with each drone 
measuring the relative position between itself and the other drones. Such 
measurements were subjected to noise during the training to make the policy
resilient during deployment in real life scenarios. This is in contrast to [8]
were the measurements were in general noise-free but did not employ any type of
inter-drone communication or measurements.


\section{Fixed wings cooperative transportation}

However, all of these works have focused on multirotor UAVs, while
comparatively little attention has been given to fixed-wing UAVs. More
recently, a new line of research has emerged exploring the same problem
using fixed-wing UAVs [3] [4] [5]. In [3], the existence of static equilibria
for a cable-suspended load transported by non-stopping flying carriers is 
investigated. The study shows that fixed-wing carriers can maintain a fixed 
pose for a cable-suspended load using at least three carriers, with some 
exceptions and degenerate cases. The work presents a theoretical feasibility 
study supported by numerical simulations in MATLAB/Simulink. However, only 
steady-state equilibrium is considered, without addressing the transient 
behavior of the system. Furthermore, the carrier reference trajectories are 
precalculated offline and executed in open loop, meaning that there is no 
reaction to disturbances, system uncertainties, or obstacles, and perfect 
coordination between the carriers is assumed. 

In[4], an algorithm for generating coordinated non-stop trajectories for 
three or more carriers holding a cable-suspended load is presented. 
The approach uses Hamiltonian cycles and graph theory to generate the 
coordinated trajectories, with the results demonstrated through numerical 
simulations in MATLAB/Simulink. However, only steady-state equilibrium is 
considered, without addressing the transient behavior of the system. 
Furthermore, the carrier reference trajectories are precalculated and there 
is no closed-loop control, meaning that the trajectories cannot react to 
disturbances or system uncertainties during operation.



In [5], a closed-loop control approach for cooperative cable-suspended 
payload manipulation using non-stopping flying carriers is presented. The 
approach considers a time-varying load trajectory, which can also include 
stopping, and consists of an outer-loop load-wrench controller, an inner-loop 
carrier-trajectory generator, and an optimization layer for the internal 
forces to enforce a strictly positive carrier velocity norm. The approach 
is experimentally validated using Crazyflie quadrotors and the Crazyswarm 
software, with a Vicon motion capture system, in addition to numerical 
simulations in MATLAB/Simulink. However, the controller is centralized,
is not deployed onboard, and assumes perfect knowledge of the system and clean,
noise-free measurements. Furthermore, the trajectories are calculated using a
centralized controller together with a Nonlinear Programming (NLP) 
solver/optimizer, resulting in a high computational cost.


The works discussed above that tackled the non-stopping carriers problems all 
relied on calculating the UAVs trajectories centrally using classical 
control techniques. While they demonstrated several effective approaches to
cooperative payload transportation using non-stopping carriers, they also 
highlighted a gap in using learning-based methods and decentralized strategies 
to tackle the same problem. To the best of the author's knowledge, this 
research direction was not explored before. Motivated 
by this gap, this work explores the use of Imitation Learning and 
Reinforcement learning as a potential research direction.
